\documentclass[11pt]{article}

\usepackage[a4paper,margin=1in]{geometry}
\usepackage{amsmath,amssymb,bm}
\usepackage{booktabs}
\usepackage{hyperref}

\hypersetup{
  colorlinks=true,
  linkcolor=blue,
  urlcolor=blue,
  citecolor=blue
}

\title{\texttt{alfak2}: Mathematical Description of the Current Method}
\author{}
\date{\today}

\newcommand{\R}{\mathbb{R}}
\newcommand{\E}{\mathbb{E}}
\newcommand{\Var}{\operatorname{Var}}
\newcommand{\Normal}{\mathcal{N}}
\newcommand{\Mult}{\operatorname{Multinomial}}
\newcommand{\DM}{\operatorname{DirichletMultinomial}}
\newcommand{\softmax}{\operatorname{softmax}}

\begin{document}
\maketitle

\section{Overview}

\texttt{alfak2} infers aneuploid karyotype fitness landscapes from sparse
two-timepoint count data. The input is a two-column count matrix
\[
  Y = \{y_{v0}, y_{v1}\}_{v \in \mathcal{O}},
\]
where $\mathcal{O}$ is the set of observed karyotypes and $y_{v0}$,
$y_{v1}$ are counts at the two timepoints. Each karyotype is represented as a
copy-number vector
\[
  x_v = (x_{v1}, \ldots, x_{vp}) \in \mathbb{Z}_{\ge 0}^{p},
\]
where $p$ is the chromosome dimension.

The current method is a two-layer model:
\[
\begin{aligned}
\text{count data}
&\longrightarrow
\text{local hierarchical Bayesian count model}
\longrightarrow
\text{local fitness anchors}\\
&\longrightarrow
\text{graph Gaussian posterior smoothing}
\longrightarrow
\text{global fitness landscape}.
\end{aligned}
\]

The first layer uses the observed counts and an explicit selection plus
missegregation forward model to estimate local fitness values and their
uncertainties. The second layer treats those local estimates as noisy anchors
and propagates information over a larger bounded karyotype graph using a sparse
Gaussian graph posterior.

\section{Notation}

\begin{table}[h]
\centering
\begin{tabular}{ll}
\toprule
Symbol & Meaning \\
\midrule
$x_v$ & Copy-number vector for node $v$ \\
$\mathcal{O}$ & Observed karyotype set \\
$\mathcal{G}=(\mathcal{V},\mathcal{E})$ & Bounded karyotype graph \\
$d(v)$ & Shortest one-step copy-number distance from observed seed nodes \\
$\beta$ & Single-chromosome missegregation probability \\
$\Delta t$ & Time interval between the two count observations \\
$T_{uv}$ & Row-stochastic transition probability from parent node $u$ to child node $v$ \\
$\eta_v$ & Initial log-abundance parameter for node $v$ \\
$\pi_{0v}, \pi_{1v}$ & Model frequencies at the two timepoints \\
$f_v$ & Native \texttt{alfak2} fitness for node $v$ \\
$z_a$ & Local fitness estimate used as a global anchor \\
$s_a^2$ & Local fitness variance estimate for anchor $a$ \\
\bottomrule
\end{tabular}
\end{table}

\section{Count Preprocessing}

If \texttt{input\_depth = "raw"}, the model uses the supplied count matrix
directly:
\[
  y'_{vt}=y_{vt}.
\]

If \texttt{input\_depth = "effective"}, each timepoint is first converted to
frequencies and then recounted at a controlled target depth:
\[
  N_t = \sum_{v \in \mathcal{O}} y_{vt},
  \qquad
  \hat p_{vt} = \frac{y_{vt}}{N_t}.
\]
The target depth is
\[
  D_t =
  \begin{cases}
  \min(N_0,N_1), & \text{mode = min},\\
  \min(N_t,D), & \text{mode = cap},\\
  D, & \text{mode = fixed}.
  \end{cases}
\]
The effective counts satisfy
\[
  y'_{vt} \approx D_t \hat p_{vt},
  \qquad
  \sum_v y'_{vt}=D_t.
\]
Rounding is performed with a largest-remainder style rule so that the total is
exactly preserved.

This step controls the influence of sequencing depth. Without it, a timepoint
with much larger depth can dominate the likelihood even when the biological
information content should be treated as comparable.

\section{Bounded Karyotype Graph}

\texttt{alfak2} builds a finite graph
$\mathcal{G}=(\mathcal{V},\mathcal{E})$ around the observed karyotypes.
Two nodes are connected by a one-step edge if they differ by one copy on one
chromosome:
\[
  x_v = x_u \pm e_c,
\]
where $e_c$ is the unit vector in chromosome coordinate $c$.

Graph expansion is bounded by copy-number limits, shell depth, and a maximum
node count:
\[
  \min_{\mathrm{cn}} \le x_{vc} \le \max_{\mathrm{cn}},
  \qquad
  d(v) \le h,
  \qquad
  |\mathcal{V}| \le M.
\]

Nodes are classified by support distance:
\[
\begin{array}{ll}
d(v)=0 & \text{directly informed},\\
d(v)=1 & \text{local borrowed},\\
d(v)=2 & \text{weakly supported},\\
d(v)>2 & \text{graph borrowed}.
\end{array}
\]

The bounded graph avoids enumerating the full copy-number lattice, whose size
grows exponentially with chromosome dimension. It keeps the inference focused
on biologically plausible states near the observed karyotypes.

\section{Missegregation Transition Kernel}

\subsection{Exact Single-Chromosome Kernel}

Let $i$ be the parent copy number and $j$ the child copy number for one
chromosome. Under the exact kernel, the transition probability is
\[
  p_{\beta}(i \to j)
  =
  \sum_{\substack{m=|i-j|\\m\equiv |i-j|\;(\mathrm{mod}\;2)}}^{i}
  \binom{i}{m}
  \beta^m(1-\beta)^{i-m}
  2^{-m}
  \binom{m}{(m+i-j)/2}.
\]
Here $m$ is the number of missegregating copies. The term
$\binom{i}{m}$ selects which copies missegregate,
$\beta^m(1-\beta)^{i-m}$ is the missegregation probability, and
$2^{-m}\binom{m}{(m+i-j)/2}$ accounts for their allocation to the two
daughter directions.

\subsection{Karyotype-Level Transition Weight}

For a full karyotype transition from node $u$ to node $v$, the exact weight is
the product over chromosomes:
\[
  q_{\beta}(u \to v)
  =
  \prod_{c=1}^{p}
  p_{\beta}(x_{uc}\to x_{vc}).
\]

The alternative \texttt{linear} kernel preserves the historical first-order
approximation:
\[
  q_{\beta}(u \to v)
  =
  \begin{cases}
  \beta\max(1,x_{uc}), & x_v=x_u\pm e_c,\\
  1, & v=u,\\
  0, & \text{otherwise}.
  \end{cases}
\]

\subsection{Row-Stochastic Transport Matrix}

The local forward model uses a row-stochastic transport matrix $T$. Let
$q_{uv}$ be the one-step edge weight and $s_u$ the self weight. Then
\[
  T_{uv}
  =
  \frac{q_{uv}}{s_u+\sum_{z:(u,z)\in\mathcal{E}}q_{uz}},
  \qquad u\ne v,
\]
and
\[
  T_{uu}
  =
  \frac{s_u}{s_u+\sum_{z:(u,z)\in\mathcal{E}}q_{uz}}.
\]
For the exact kernel, $s_u=q_\beta(u\to u)$. For the linear kernel, $s_u=1$.

\section{Local Hierarchical Bayesian Count Model}

The local model estimates initial frequencies and fitness values on a local
karyotype graph $\mathcal{G}_{\mathrm{local}}$.

\subsection{Initial Frequency Parameterization}

Each node $v$ has an initial log-abundance parameter $\eta_v$. Initial
frequencies are obtained by a softmax transform:
\[
  \pi_{0v}
  =
  \frac{\exp(\eta_v)}
       {\sum_{u\in\mathcal{V}}\exp(\eta_u)}.
\]

\subsection{Selection and Missegregation Forward Model}

Fitness acts through exponential growth between the two timepoints:
\[
  g_v = \pi_{0v}\exp(\Delta t\,f_v).
\]
After growth, mass is transported through the missegregation matrix:
\[
  \tilde\pi_{1v}
  =
  \sum_{u\in\mathcal{V}} T_{uv}g_u.
\]
The result is normalized to a frequency vector:
\[
  \pi_{1v}
  =
  \frac{\tilde\pi_{1v}}
       {\sum_{z\in\mathcal{V}}\tilde\pi_{1z}}.
\]

Thus the model explicitly represents
\[
  \text{initial composition}
  \quad+\quad
  \text{fitness-driven growth}
  \quad+\quad
  \text{missegregation transport}.
\]

\subsection{Observation Model}

For raw counts, the default observation model is multinomial:
\[
  y_{\cdot 0}\mid \pi_0
  \sim
  \Mult(N_0,\pi_0),
  \qquad
  y_{\cdot 1}\mid \pi_1
  \sim
  \Mult(N_1,\pi_1),
\]
where $N_t=\sum_v y_{vt}$. Up to constants independent of the parameters,
\[
  \ell_{\mathrm{mult}}
  =
  \sum_v y_{v0}\log \pi_{0v}
  +
  \sum_v y_{v1}\log \pi_{1v}.
\]

The optional Dirichlet-multinomial observation model is
\[
  y_{\cdot t}\mid \pi_t,\phi
  \sim
  \DM(N_t,\pi_t,\phi),
\]
with log-likelihood
\[
  \ell_{\mathrm{DM}}(y,\pi,\phi)
  =
  \log\Gamma(\phi)
  -
  \log\Gamma(N+\phi)
  +
  \sum_v
  \left[
    \log\Gamma(y_v+\phi\pi_v)
    -
    \log\Gamma(\phi\pi_v)
  \right].
\]
This model allows overdispersion relative to the multinomial likelihood.

\subsection{Prior for Initial Abundance}

Initial log-abundance has a node-specific Gaussian prior:
\[
  \eta_v \sim \Normal(m_{\eta v},s_{\eta v}^2).
\]
Directly observed or otherwise effectively counted nodes use a diffuse prior:
\[
  m_{\eta v}=0,\qquad s_{\eta v}=5.
\]
Unobserved borrowed nodes use a smaller prior mean:
\[
  m_{\eta v}
  =
  -6 - 0.75\max\{0,d(v)-1\},
  \qquad
  s_{\eta v}=1.5.
\]
This encodes the conservative assumption that unobserved nearby karyotypes may
exist but should have small initial abundance unless the data support them.

\subsection{Hierarchical Fitness Prior}

For directly informed nodes, fitness is centered at zero:
\[
  f_v\mid \sigma_{\mathrm{anchor}}
  \sim
  \Normal(0,\sigma_{\mathrm{anchor}}^2),
  \qquad d(v)=0.
\]

Borrowed nodes receive information from parent nodes closer to the observed
set. Each parent edge $(u,v)$ has a context
\[
  c(u,v)
  =
  \{\text{gain/loss},\;\text{chromosome index},\;\text{parent ploidy band}\}.
\]
Context effects follow a hierarchical Gaussian prior:
\[
  \delta_c\mid \mu_{g(c)},\tau_{g(c)}
  \sim
  \Normal(\mu_{g(c)},\tau_{g(c)}^2),
\]
\[
  \mu_g\sim \Normal(0,0.5^2).
\]
The scale parameters are positive and are given half-normal type priors:
\[
  \sigma_{\mathrm{neighbor}}>0,\qquad
  \sigma_{\mathrm{anchor}}>0,\qquad
  \tau_g>0.
\]

For a borrowed node $v$, the prior mean is the transition-weighted average of
parent fitness plus context effect:
\[
  m_{fv}
  =
  \frac{
    \sum_{u\in P(v)}
    w_{uv}\{f_u+\delta_{c(u,v)}\}
  }{
    \sum_{u\in P(v)}w_{uv}
  }.
\]
The borrowed-node prior is
\[
  f_v\mid \{f_u:u\in P(v)\},\delta,\sigma_{\mathrm{neighbor}}
  \sim
  \Normal(m_{fv},\sigma_{\mathrm{neighbor}}^2),
  \qquad d(v)>0.
\]
When $d(v)\ge 2$, the local model inflates this prior standard deviation by
$1.75$:
\[
  f_v \sim \Normal(m_{fv},(1.75\sigma_{\mathrm{neighbor}})^2).
\]
This expresses weaker local information farther from observed nodes.

\subsection{Local Posterior Objective}

The local layer estimates the maximum a posteriori parameter vector
\[
  \theta=(\eta,f,\delta,\mu,\sigma_{\mathrm{neighbor}},
  \sigma_{\mathrm{anchor}},\tau).
\]
It solves
\[
  \hat\theta
  =
  \arg\max_{\theta}
  \left\{
    \ell(Y\mid \eta,f,\beta,\Delta t)
    +
    \log p(\eta)
    +
    \log p(f\mid \delta)
    +
    \log p(\delta,\mu,\sigma,\tau)
  \right\}.
\]
Equivalently, the code minimizes the negative log posterior
\[
  \mathcal{J}_{\mathrm{local}}(\theta)
  =
  -\ell(Y\mid \eta,f,\beta,\Delta t)
  -
  \log p(\eta)
  -
  \log p(f\mid \delta)
  -
  \log p(\delta,\mu,\sigma,\tau).
\]

This objective is implemented with Template Model Builder (TMB). TMB is not the
statistical algorithm; it is the numerical backend used to compute derivatives,
optimize the MAP objective, and obtain approximate marginal standard errors
through \texttt{sdreport}.

The local layer returns, for each local graph node,
\[
  z_v=\hat f_v,\qquad s_v=\widehat{\mathrm{sd}}(f_v).
\]
These become noisy anchors for the global graph layer.

\section{Global Graph Gaussian Posterior}

The global layer does not revisit the raw counts. Instead, it treats local
fitness estimates as noisy observations and infers the full fitness vector
\[
  \bm f=(f_1,\ldots,f_n)^\top
\]
on a larger graph $\mathcal{G}_{\mathrm{global}}$.

\subsection{Anchor Observation Model}

Let $\mathcal{A}$ be the retained anchor set. For each anchor $a$,
\[
  z_a\mid f_a
  \sim
  \Normal(f_a,R_a),
\]
where
\[
  R_a = s_a^2\,m_a + \sigma_{\mathrm{obs}}^2.
\]
The multiplier $m_a$ combines covariance-status and count-based inflation:
\[
  m_a
  =
  m^{\mathrm{cov}}_a
  m^{\mathrm{count}}_a.
\]
If count-based downweighting is enabled,
\[
  m^{\mathrm{count}}_a
  =
  \max\left\{
    1,
    \left(\frac{C_{\mathrm{ref}}}{C_a}\right)^\alpha
  \right\}.
\]
Low-count or unreliable local anchors therefore receive larger observation
variance and have less influence on the global posterior.

\subsection{Graph Laplacian Smoothing}

The first global prior term is graph edge smoothing:
\[
  \mathcal{P}_{\mathrm{edge}}(\bm f)
  =
  \sum_{(i,j)\in\mathcal{E}}
  w_{ij}(f_i-f_j)^2.
\]
It enters the posterior with strength $\lambda_l$:
\[
  \lambda_l\mathcal{P}_{\mathrm{edge}}(\bm f).
\]
This term encourages neighboring karyotypes connected by plausible one-step
copy-number changes to have similar fitness unless anchor data support a sharp
difference.

\subsection{Ordered Copy-Number Curvature Penalty}

The second prior term penalizes second differences along a chromosome
copy-number axis. For any triplet $x-e_c$, $x$, $x+e_c$ present in the graph,
define
\[
  \Delta_c^2 f(x)
  =
  f(x-e_c)-2f(x)+f(x+e_c).
\]
The curvature penalty is
\[
  \mathcal{P}_{\mathrm{curv}}(\bm f)
  =
  \sum_{x,c}
  \left\{\Delta_c^2 f(x)\right\}^2.
\]
This discourages unsupported zig-zag behavior along ordered copy-number
coordinates.

\subsection{Epistasis Penalty}

The third prior term controls second-order interaction between two chromosome
directions. For any quadruplet $x$, $x+e_c$, $x+e_d$, and $x+e_c+e_d$ present
in the graph,
\[
  \Delta_{cd} f(x)
  =
  f(x)-f(x+e_c)-f(x+e_d)+f(x+e_c+e_d).
\]
The epistasis penalty is
\[
  \mathcal{P}_{\mathrm{epi}}(\bm f)
  =
  \frac{1}{2}
  \sum_{x,c<d}
  \left\{\Delta_{cd}f(x)\right\}^2.
\]
This does not forbid epistasis. It penalizes strong interaction effects when
they are weakly supported by anchors, reducing unstable high-order oscillations
in sparse regions.

\subsection{Gaussian Posterior}

The global negative log posterior is
\[
\begin{aligned}
  \mathcal{J}_{\mathrm{global}}(\bm f)
  =
  \frac{1}{2}
  \bigg[
  &\lambda_l\mathcal{P}_{\mathrm{edge}}(\bm f)
  +
  \lambda_e\mathcal{P}_{\mathrm{curv}}(\bm f)
  +
  \lambda_e\mathcal{P}_{\mathrm{epi}}(\bm f)\\
  &+
  \sum_{a\in\mathcal{A}}
  \frac{(f_a-z_a)^2}{R_a}
  +
  \epsilon\sum_{v\in\mathcal{V}}f_v^2
  \bigg].
\end{aligned}
\]
All terms are quadratic, so the posterior is Gaussian:
\[
  p(\bm f\mid \bm z)
  =
  \Normal(\bm \mu,\bm Q^{-1}).
\]
The precision matrix can be written as
\[
  \bm Q
  =
  \lambda_l\bm L
  +
  \lambda_e\bm K
  +
  \bm A^\top\bm R^{-1}\bm A
  +
  \epsilon\bm I.
\]
Here $\bm L$ is the weighted graph Laplacian, $\bm K$ is the sparse precision
contribution from the ordered copy-number curvature and epistasis penalties,
$\bm A$ selects anchor nodes from the full graph, and $\bm R$ is the diagonal
anchor variance matrix.

The right-hand side is
\[
  \bm b
  =
  \bm A^\top\bm R^{-1}\bm z.
\]
The posterior mean solves
\[
  \bm Q\bm\mu=\bm b,
  \qquad
  \bm \mu=\bm Q^{-1}\bm b.
\]
The marginal posterior variance for node $v$ is
\[
  \Var(f_v\mid \bm z)
  =
  (\bm Q^{-1})_{vv}.
\]

The implementation assembles $\bm Q$ as a sparse matrix and uses RcppEigen
sparse linear algebra to solve for $\bm \mu$ and the diagonal of
$\bm Q^{-1}$.

\subsection{Hyperparameter Selection}

The global layer has three main hyperparameters:
\[
  \lambda_l,\qquad \lambda_e,\qquad \sigma_{\mathrm{obs}}.
\]
If each grid contains a single value, these are treated as fixed. If multiple
values are supplied, \texttt{alfak2} uses leave-one-anchor-out cross-validation.
For each anchor $a$, the model removes $a$, predicts its fitness from the
remaining anchors, and obtains $\mu_{-a,a}$. The score is
\[
  \operatorname{CV}(\lambda_l,\lambda_e,\sigma_{\mathrm{obs}})
  =
  \frac{1}{|\mathcal{A}_{\mathrm{eval}}|}
  \sum_{a\in\mathcal{A}_{\mathrm{eval}}}
  \frac{(z_a-\mu_{-a,a})^2}{R_a}.
\]
The selected values are
\[
  (\lambda_l^\star,\lambda_e^\star,\sigma_{\mathrm{obs}}^\star)
  =
  \arg\min
  \operatorname{CV}(\lambda_l,\lambda_e,\sigma_{\mathrm{obs}}).
\]

\section{Output}

The local layer reports
\[
  \hat f_v^{\mathrm{local}},
  \qquad
  \widehat{\operatorname{sd}}(f_v^{\mathrm{local}}),
  \qquad
  \hat\pi_{0v},
  \qquad
  \hat\pi_{1v}.
\]
The global layer reports
\[
  \hat f_v^{\mathrm{global}}=\mu_v,
  \qquad
  \widehat{\operatorname{sd}}(f_v^{\mathrm{global}})
  =
  \sqrt{(\bm Q^{-1})_{vv}}.
\]
Approximate normal intervals are
\[
  \hat f_v \pm 1.9599639845\,\widehat{\operatorname{sd}}(f_v).
\]

The fit also reports support tier, support distance, anchor diagnostics,
selected hyperparameters, and cross-validation status. Nodes far from anchors
or with large posterior standard deviation can be labeled as graph-borrowed or
prior-dominated, indicating that their estimates are driven mostly by prior
structure rather than direct data.

\section{Optional \texttt{alfakR}-Scale Transformation}

Native \texttt{alfak2} fitness is on a relative scale. If
\texttt{alfakR\_scale = TRUE}, the package appends legacy-scale columns for
compatibility with \texttt{alfakR}. Let
\[
  g_0=\frac{\log(n_b/n_0)}{\Delta t}
\]
and define the efflux viability approximation
\[
  V(x)=2(1-\beta)^{\sum_c x_c}-1.
\]

Without efflux correction,
\[
  f^{\mathrm{alfakR}}(x)
  =
  f^{\mathrm{alfak2}}(x)
  +
  g_0
  -
  \sum_x w_x f^{\mathrm{alfak2}}(x).
\]
With efflux correction,
\[
  o
  =
  \frac{
    \sum_x w_x f^{\mathrm{alfak2}}(x)/V(x)-g_0
  }{
    \sum_x w_x/V(x)
  },
\]
\[
  f^{\mathrm{alfakR}}(x)
  =
  \frac{f^{\mathrm{alfak2}}(x)-o}{V(x)}.
\]
This is a post-processing scale conversion. It does not change the native
\texttt{alfak2} posterior estimates.

\section{Why This Method Is Useful}

\begin{enumerate}
  \item \textbf{It uses count likelihoods directly.}
  The local layer models counts with multinomial or Dirichlet-multinomial
  likelihoods rather than only comparing observed frequency changes.

  \item \textbf{It separates selection from missegregation.}
  The forward model
  \[
    \pi_0 \to \pi_0\exp(\Delta t f)
    \to T^\top\{\pi_0\exp(\Delta t f)\}
  \]
  represents both fitness-driven growth and missegregation transport.

  \item \textbf{It is designed for sparse two-timepoint data.}
  Nearby unobserved nodes borrow information through parent fitness and context
  effects, while their uncertainty increases with weaker support.

  \item \textbf{It avoids full lattice enumeration.}
  Bounded graph expansion makes high-dimensional karyotype spaces
  computationally manageable.

  \item \textbf{It uses structured graph smoothing.}
  The global layer uses graph Laplacian smoothing, ordered copy-number
  curvature penalties, and epistasis penalties that reflect the structure of
  copy-number landscapes.

  \item \textbf{It preserves uncertainty.}
  The local TMB layer provides anchor uncertainty, and the global Gaussian
  posterior provides posterior standard deviations for all graph nodes.

  \item \textbf{It provides interpretable support tiers.}
  Each node is annotated by support distance and support tier, helping users
  distinguish directly informed estimates from graph-borrowed estimates.
\end{enumerate}

\section{Summary}

The current \texttt{alfak2} method can be summarized as
\[
\begin{aligned}
\text{Local:}\quad
&\text{counts}
+ \text{evolution dynamics}
+ \text{neighbor priors}\\
&\longrightarrow
\text{fitness anchors};
\\[4pt]
\text{Global:}\quad
&\text{anchor observations}
+ \text{graph smoothing}
+ \text{CN penalties}\\
&\longrightarrow
\text{graph Gaussian posterior}.
\end{aligned}
\]

\end{document}
